% Part: normal-modal-logic
% Chapter: syntax-and-semantics
% Section: relational-models

\documentclass[../../../include/open-logic-section]{subfiles}

\begin{document}

\olfileid{nml}{syn}{rel}

\olsection{النماذج العلائقية}

مدار دلالة المنطقات الجهية النظامية على \emph{النموذج العلائقي}.
نأخذ مجموعة من العوالم، ونجعل بينها علاقة ثنائية تسمى علاقة الوصول،
ثم نعيّن أي !!{propositional variable}s تصدق في كل عالم. فبهذه
الأجزاء الثلاثة يتحدد النموذج.

\begin{defn}
  \emph{نموذج} اللغة الموجهية الأساسية ثلاثية~$\mModel{M}
  = \tuple{{\OLId{latin-looped}{W}}, {\OLId{latin-looped}{R}}, {\OLId{latin-looped}{V}}}$، حيث
  \begin{enumerate}
  \item ${\OLId{latin-looped}{W}}$ مجموعة غير خالية من ``العوالم''،
  \item ${\OLId{latin-looped}{R}}$ علاقة وصول ثنائية على~${\OLId{latin-looped}{W}}$، و
  \item ${\OLId{latin-looped}{V}}$ دالة تعين لكل !!{propositional variable}~${\OLId{latin-ordinary}{p}}$ مجموعة
    ${\OLId{latin-looped}{V}}({\OLId{latin-ordinary}{p}})$ من العوالم الممكنة.
  \end{enumerate}
  وعندما يصدق~$Rww'$، نقول إن~${\OLId{latin-ordinary}{w}}'$ \emph{يمكن الوصول إليه من}~${\OLId{latin-ordinary}{w}}$.
  وعندما يكون~${\OLId{latin-ordinary}{w}} \in {\OLId{latin-looped}{V}}({\OLId{latin-ordinary}{p}})$، نقول إن~${\OLId{latin-ordinary}{p}}$ \emph{صادق في}~${\OLId{latin-ordinary}{w}}$.
\end{defn}

ومن فائدة الدلالة العلائقية أن نموذجها يُرى في رسم يسير، كما
في~\olref{fig:simple}. فالعوالم عقد، ووصول العالم~${\OLId{latin-ordinary}{w}}'$ من~${\OLId{latin-ordinary}{w}}$
يمثله سهم يخرج من~${\OLId{latin-ordinary}{w}}$ وينتهي إلى~${\OLId{latin-ordinary}{w}}'$. ونضع على العقدة
$\mTrue{{\OLId{latin-ordinary}{p}}}$ حين يكون~${\OLId{latin-ordinary}{w}} \in {\OLId{latin-looped}{V}}({\OLId{latin-ordinary}{p}})$، ونضع~\mFalse{p} فيما عدا ذلك.
وعليه فإن~\olref{fig:simple} يرسم النموذج الذي تعيّنه البيانات
${\OLId{latin-looped}{W}} = \{{\OLId{latin-ordinary}{w}}_{\OLMathNumeral{1}}, {\OLId{latin-ordinary}{w}}_{\OLMathNumeral{2}}, {\OLId{latin-ordinary}{w}}_{\OLMathNumeral{3}}\}$، و${\OLId{latin-looped}{R}} = \{\tuple{{\OLId{latin-ordinary}{w}}_{\OLMathNumeral{1}}, {\OLId{latin-ordinary}{w}}_{\OLMathNumeral{2}}},\tuple{{\OLId{latin-ordinary}{w}}_{\OLMathNumeral{1}},
  {\OLId{latin-ordinary}{w}}_{\OLMathNumeral{3}}}\}$، و${\OLId{latin-looped}{V}}({\OLId{latin-ordinary}{p}}) = \{{\OLId{latin-ordinary}{w}}_{\OLMathNumeral{1}}, {\OLId{latin-ordinary}{w}}_{\OLMathNumeral{2}}\}$، و${\OLId{latin-looped}{V}}({\OLId{latin-ordinary}{q}}) = \{{\OLId{latin-ordinary}{w}}_{\OLMathNumeral{2}}\}$.

\begin{figure}
  \begin{center}
    \begin{tikzpicture}[modal]
      \node[world] (w1) [label={[align=right]right:\mTrue{p}\\ \mFalse{q}}]{${\OLId{latin-ordinary}{w}}_{\OLMathNumeral{1}}$};
      \node[world] (w2) [label={[align=right]right:\mTrue{p}\\ \mTrue{q}},
        above right=of w1]{${\OLId{latin-ordinary}{w}}_{\OLMathNumeral{2}}$};
      \node[world] (w3) [label={[align=right]right:\mFalse{p}\\ \mFalse{q}},
        below right=of w1] {${\OLId{latin-ordinary}{w}}_{\OLMathNumeral{3}}$};
      \draw[->] (w1) to (w2);
      \draw[->] (w1) to (w3);
    \end{tikzpicture}
  \end{center}
  \caption{نموذج بسيط.}
  \ollabel{fig:simple}
\end{figure}

\end{document}
